% titel: Blatt 5, A1: Eigenschaften von Funktionen
% Teil II: der durchgerechnete Weg.

\subsection*{Gegeben und gesucht}
Drei Funktionen auf Gerade/Ungerade und Monotonie untersuchen (a), einen
Satz über die Umkehrfunktion einer streng monoton steigenden Funktion
begründen (b), und eine bijektive Funktion $f\colon[0,1]\to[0,1]$ finden,
die nicht streng monoton steigend ist (c).

\subsection*{a) Gerade/Ungerade und Monotonie}
\emph{Gedankenschritt:} Symmetrie prüft man durch Ersetzen von $x$ durch
$-x$; Monotonie durch Vergleich zweier Argumente $x_1<x_2$. Beide Prüfungen
sind unabhängig voneinander durchzuführen.

\begin{itemize}
  \item $f_1(x)=\sqrt{1-x^2}$ auf $[-1,1]$. Einsetzen:
        $f_1(-x)=\sqrt{1-(-x)^2}=\sqrt{1-x^2}=f_1(x)$, also \textbf{gerade}.
        Monotonie: Auf $[0,1]$ fällt das Argument $1-x^2$ mit wachsendem $x$;
        die (steigende) Wurzel überträgt das Fallen, also ist $f_1$ dort
        streng fallend, auf $[-1,0]$ umgekehrt streng steigend. Wegen des
        Symmetrieknicks bei $0$ ist $f_1$ auf \emph{ganz} $[-1,1]$ weder
        streng steigend noch streng fallend.
  \item $f_2(x)=10-10x$. Einsetzen: $f_2(-x)=10+10x$, weder $f_2(x)$ noch
        $-f_2(x)$ -- $f_2$ ist weder gerade noch ungerade. Monotonie: Aus
        $x_1<x_2$ folgt $-x_1>-x_2$, also $10-10x_1>10-10x_2$, d.h.
        $f_2(x_1)>f_2(x_2)$: \textbf{streng monoton fallend}. (Der Graph ist
        eine Gerade mit negativer Steigung $-10$; Steigung und Monotonie
        stimmen überein.)
  \item $f_3(x)=x^2\sin x$. Einsetzen:
        $f_3(-x)=(-x)^2\sin(-x)=x^2(-\sin x)=-x^2\sin x=-f_3(x)$, also
        \textbf{ungerade}. Der Gedanke dahinter: $x^2$ ist gerade, $\sin x$
        ungerade, und \emph{gerade mal ungerade gleich ungerade}. Monotonie: Wegen
        des oszillierenden $\sin$ mit wachsender Amplitude $x^2$ kann $f_3$
        auf $\mathbb R$ nicht monoton sein.
\end{itemize}

\subsection*{b) Umkehrfunktion ist streng monoton steigend}
\emph{Warum die Kontraposition?} Wir zeigen, dass die Umkehrfunktion $g$
steigt. Die einzige Information über ihre Wirkung ist $f(g(u))=u$. Direkt
kann man sie nicht zerlegen, also nimmt man an, $g$ wäre \emph{nicht}
steigend, und leitet daraus einen Widerspruch ab (Kontraposition).

Seien $u<w$ aus dem Wertebereich. Angenommen $g(u)\ge g(w)$. Weil $f$
streng monoton steigend ist, gilt für die neuen Argumente
$f(g(u))\ge f(g(w))$. Nach $f\circ g=\mathrm{id}$ ist das aber genau
$u\ge w$ -- im Widerspruch zu $u<w$. Die Annahme war falsch, also ist
$g(u)<g(w)$; $g$ ist streng monoton steigend.

\subsection*{c) Bijektiv, aber nicht streng monoton steigend}
$g(x)=1-x$ bildet $[0,1]$ bijektiv auf $[0,1]$ ab (sie ist ihre eigene
Umkehrung: $g(g(x))=1-(1-x)=x$) und ist wegen der negativen Steigung $-1$
streng monoton fallend -- also sicher nicht streng monoton steigend. Das
Gegenbeispiel zeigt: \emph{Bijektivität} und \emph{Monotonie} sind
unabhängige Eigenschaften.

\begin{ergebnis}
\begin{itemize}
  \item $f_1$ gerade; $f_2$ streng monoton fallend; $f_3$ ungerade.
  \item Die Umkehrfunktion einer streng monoton steigenden surjektiven
        Funktion ist streng monoton steigend.
  \item Ja: $g(x)=1-x$ ist bijektiv auf $[0,1]$, aber nicht monoton
        steigend.
\end{itemize}
\end{ergebnis}

\begin{kontrolle}
\textbf{Probe:} $g(g(x))=1-(1-x)=x$, also ist $g$ die Identität auf
$[0,1]$ -- $g$ ist ihre eigene Umkehrung und damit bijektiv. Für
$x_1<x_2$ ist $g(x_1)=1-x_1>1-x_2=g(x_2)$, also streng monoton fallend.
\checkmark

\textbf{Grenzfall (a):} $f_1$ hat bei $x=0$ ein Maximum $1$ und gilt
$f_1(\pm1)=0$. Eine gerade Funktion muss an beiden Enden denselben Wert
annehmen ($f_1(-1)=f_1(1)=0$); das verträgt sich nicht mit strenger
Monotonie auf ganz $[-1,1]$. \checkmark
\end{kontrolle}

\begin{merksatz}
Gerade/ungerade prüft man durch Einsetzen von $-x$; die Symmetrie des
Produkts ist das \emph{Produkt der Symmetrien}. Streng monotone Funktionen
sind injektiv und damit umkehrbar; die Umkehrfunktion erbt die
Monotonierichtung (Beweis über Kontraposition). Bijektivität erzwingt
keine Monotonie (Gegenbeispiel $x\mapsto 1-x$).
\end{merksatz}
